\section{Geo Dynamics Conventions}

\subsection{Purpose}

The Geo path is a clean restart of the collision dynamics code.  Its purpose is
to build a Python implementation whose structure can be transferred to GLSL
with minimal translation.

\subsection{File Boundaries}

\texttt{GeoDynamics.py} is the GLSL-portable dynamics layer.  It should contain
only functions that correspond directly to shader-side dynamics functions.

\texttt{GeoBase.py} is the Python harness.  It may load configuration files,
create Python data structures, build particle arrays, set push constants,
iterate particles, and call the Geo dynamics functions.

\texttt{Reporting.py} is reserved for reporting and diagnostic output.  It
should not own dynamics decisions.

\subsection{Naming Convention}

Functions intended to map to GLSL should use GLSL-like names and signatures
where practical.  Current examples are:

\begin{verbatim}
GeoAddParticleContact(SourceID, TargetID)
GeoProcessCollisions(SourceID)
GeoMoveParticle(SourceID, positionBuffer, dt)
isParticleContact(Frame, SourceID, TargetID, positionBuffer)
\end{verbatim}

\texttt{GeoDynamics.py} should avoid Python-only helper names unless the helper
is also expected to exist in GLSL.

Before adding or renaming any Geo field, buffer, function, or constant, check
the GLSL source files first.  The GLSL names are the naming authority for the
Geo path:

\begin{verbatim}
C:\_DJ\gPCD\vulkan\sim\common\*.glsl
C:\_DJ\gPCD\vulkan\sim\BoundarySpheresMotion\BoundarySpheresMotionP.comp
\end{verbatim}

Python names should stay coordinated with those files whenever possible.  If a
Python-only harness name is needed, it belongs in \texttt{GeoBase.py},
\texttt{GeoRunner.py}, or \texttt{Reporting.py}, not in
\texttt{GeoDynamics.py}.

\subsection{Data Convention}

Particle records use field access:

\begin{verbatim}
Particle.vx
Particle.radius
\end{verbatim}

rather than dictionary access:

\begin{verbatim}
Particle["vx"]
\end{verbatim}

The Python structures should be shaped to resemble GLSL buffers and push
constants.  Python-only configuration parsing must not leak into
\texttt{GeoDynamics.py}.

When adding particle position, velocity, radius, collision-list, push-constant,
or error fields, first look for the corresponding GLSL structure or function
name in the common shader files.  For example, position double buffering should
prefer names coordinated with the GLSL particle buffers and
\texttt{ShaderFlags.positionBuffer}, rather than inventing unrelated Python-only
names.

\subsection{Overlap Area Convention}

Overlap area is a geometric contact property.  For a particle-particle contact,
the source-side and target-side overlap areas are the same geometric area, even
when the two particles have different velocities.  Therefore, once the overlap
area is computed for a source-target pair, the target particle's overlap area is
known from the same calculation.

Velocity may affect momentum transfer, phase, or future position, but it does
not make the instantaneous geometric overlap area different for the two
particles in the same contact.

\subsection{Instantaneous Geometric Goal}

The long-term Geo dynamics target is an instantaneous, purely geometric
collision model.

Instantaneous means the collision response for a frame is computed from the
current frame's particle state only:

\begin{itemize}
    \item current positions;
    \item current velocities;
    \item current radii and masses;
    \item current overlap geometry;
    \item current source-owned contact list for this frame.
\end{itemize}

Purely geometric means the response should be derived from contact geometry,
such as center distance, normal direction, penetration depth, overlap area, and
the distribution of overlap area across a source particle's simultaneous
contacts.  Energy calculations and energy-named dynamics variables should be
reserved for diagnostics only.

The model should not require persistent collision-episode memory such as:

\begin{itemize}
    \item stored contact momentum reservoirs;
    \item immutable first-contact velocity;
    \item persistent compression/rebound phase state;
    \item remembered alpha-zero values;
    \item global wall or contact reservoirs shared by particles.
\end{itemize}

The only persistent particle state should be ordinary particle state, such as
position, velocity, radius, mass, and explicitly approved per-particle fields.
The current frame may build temporary contact information, but that temporary
contact information should be safe to rebuild on the next frame.

\subsection{Instantaneous Frame Snapshot Rule}

The collision response for a frame must not depend on velocity changes induced
earlier in the same frame by other contacts.  Those changes are schedule
artifacts, not instantaneous state.  If two executions see the same frame-start
positions, velocities, radii, masses, contact geometry, and approved
source-owned carried state, they must compute the same dynamics result
independent of source invocation order.

This means collision-induced velocity changes are not a separate dynamics
input.  A frame is treated as one possible instantaneous state that could have
come from many prior histories.  The response may use:

\begin{itemize}[leftmargin=*]
    \item frame-start positions for contact geometry;
    \item frame-start velocities for relative normal velocity and available
    momentum;
    \item current radii, masses, and particle-owned stiffness;
    \item current-frame contact sets and overlap-area weights;
    \item explicitly approved source-owned carried state, such as the
    contact-slot internal normal momentum value.
\end{itemize}

The response must not use:

\begin{itemize}[leftmargin=*]
    \item a target velocity already modified by another source invocation in
    the same frame;
    \item an intra-frame contact response as if it were collision history;
    \item source-order side effects to decide compression, return, or release.
\end{itemize}

Python may satisfy this rule with a frame-start velocity snapshot such as
\texttt{VelRadFrame}.  GLSL must provide an equivalent frame-start read path
before multi-contact behavior can be considered shader-equivalent.

\subsection{Approved Source-Owned Internal Momentum Exception}

The Geo path may carry one source-owned internal normal momentum value per
active contact slot.  This is an explicit exception to the fully stateless
geometric target because the current overlap-area impulse needs a ledger to
avoid releasing more momentum than it removed during compression.

The exception is constrained:

\begin{itemize}[leftmargin=*]
    \item the value is owned by the source particle's contact slot;
    \item it is not a shared pair reservoir;
    \item it must not write the target particle;
    \item it may be used only to cap source-owned rebound/release impulse;
    \item it must remain representable in the GLSL \texttt{ContactState}
    structure.
\end{itemize}

In the current GLSL-shaped layout, the carried value is:

\begin{verbatim}
P[SourceID].contacts[slot].aux.z
\end{verbatim}

This field stores source-owned internal normal momentum.  It exists to prevent
symmetric velocity drift while preserving the source-owned shader execution
model.

\subsection{Model-Preservation Warning}

Before editing code, any proposed change that would move the Geo path away from
the instantaneous, purely geometric target must be called out explicitly.

Examples of changes that require warning before implementation:

\begin{itemize}
    \item adding persistent per-contact memory for dynamics;
    \item adding a reservoir whose value survives after the frame;
    \item using first-contact velocity as a long-lived dynamics input;
    \item using phase history to decide future dynamics;
    \item adding global writable memory shared by all particles;
    \item adding a wall-specific dynamics path instead of expressing the wall as
    contact geometry.
\end{itemize}

This warning does not forbid such changes forever, but it makes the model cost
visible before implementation.

\subsection{Documenting Problem Evolution}

When a dynamics rule changes, the reason for the change should be recorded in
the problem-evolution document.  That document is the design history for the
Geo path.  It should explain how the problem was observed, which fixes were
rejected, which rule or exception was accepted, and what result verified the
change.

New entries should follow this pattern:

\begin{itemize}[leftmargin=*]
    \item identify the observed failure by test or symptom;
    \item include the commit hash associated with the change;
    \item describe rejected fixes when they would violate the GLSL execution
    model;
    \item state any convention, rule, or exception that was added;
    \item describe the implemented algorithm in equations or direct procedural
    terms;
    \item record the verification result that shows the issue was resolved.
\end{itemize}

The purpose is not only to document the current code.  It is also to preserve
the assumptions behind that code so future work can deliberately keep, revise,
or remove those assumptions.

\subsection{Control Flow Shape}

The Python control-flow harness should match the compute shader shape:

\begin{verbatim}
for each SourceID:
    build source-owned contact list
    GeoProcessCollisions(SourceID)
    GeoMoveParticle(SourceID, positionBuffer, dt)
\end{verbatim}

The current Python harness may use a direct all-pairs scan while the shader uses
cell/corner occupancy lookup.  That difference belongs in \texttt{GeoBase.py},
not \texttt{GeoDynamics.py}.

\subsection{Sequential Stage And Narrow Function Rule}

Geo frame processing should use a sequential, data-oriented stage chain.  The
stage chain owns the visible decisions about what kind of work happens next.
The narrow functions called by the chain should each perform one kind of
activity.

For example, a frame-level chain may decide:

\begin{verbatim}
build frame snapshot
build particle contact lists
build wall contact lists
resolve particle contacts
resolve wall contacts
move particles
report diagnostics
\end{verbatim}

The functions inside those stages should stay narrow:

\begin{itemize}[leftmargin=*]
    \item \texttt{isParticleContact(...)} answers only whether two particles
    are touching in the current snapshot;
    \item \texttt{isWallContact(...)} answers only whether a particle touches a
    wall or wall particle in the current snapshot;
    \item \texttt{GeoAddParticleContact(...)} records a particle contact only;
    \item \texttt{GeoAddWallContact(...)} records a wall contact only;
    \item geometry functions compute geometry only;
    \item impulse functions compute impulse values only;
    \item apply functions write the source-owned state only.
\end{itemize}

The narrower the split, the more situations can share the same downstream
dynamics.  Particle-particle contacts, wall contacts, and wall ghost-particle
contacts should differ only where they truly differ.  Usually that is contact
determination and contact creation:

\begin{verbatim}
isParticleContact(...)
isWallContact(...)
GeoAddParticleContact(...)
GeoAddWallContact(...)
\end{verbatim}

After a contact has been expressed as source-owned contact state, later stages
should be shared whenever possible:

\begin{verbatim}
GeoContactGeometry(...)
GeoContactContext(...)
GeoContactWeight(...)
GeoCompressionImpulse(...)
GeoReturnImpulse(...)
GeoApplySourceImpulse(...)
\end{verbatim}

This prevents the code from growing separate particle and wall dynamics paths
when both contacts can use the same response rule.  If a situation needs a
different narrow function, add that narrow function at the point of real
difference instead of cloning a broad stage.

Do not add an umbrella stage when it only hides useful frame flow.  For
example, this is less clear:

\begin{verbatim}
GeoBuildContactLists()
\end{verbatim}

than leaving the meaningful decisions visible in the main chain:

\begin{verbatim}
GeoBuildParticleContactLists()
GeoBuildWallContactLists()
\end{verbatim}

Umbrella functions are allowed only when they remove real duplication or hide
mechanics that are not relevant to understanding the frame sequence.

This style makes the program easier to step through, easier to debug, and
easier to learn.  If a contact list is wrong, the defect should be local to the
contact-detection or contact-recording stage.  If an impulse is wrong, the
defect should be local to the context, weighting, or impulse stage.  Decisions
about which narrow functions are called should remain visible in the stage
chain instead of being hidden inside broad helper functions.

This rule applies to the Geo dynamics path and its frame harness.  Unrelated
Python utilities do not need to be rewritten to match this style.

\subsection{No Pairwise Write-Back Shortcut}

The Geo Python implementation must not use a pairwise dynamics shortcut that
cannot be exported to the current GLSL execution model.

The active shader shape is source-owned: one invocation owns one source
particle.  Therefore a Geo dynamics function may write only the current source
particle's state.  It must not solve a pair by directly updating both
particles:

\begin{verbatim}
solve_pair(source, target):
    P[source].velocity = ...
    P[target].velocity = ...   // forbidden in GeoDynamics.py
\end{verbatim}

This rule applies even when a pairwise update would be physically cleaner in
Python.  Python is not allowed to become a separate reference model that fixes
collisions in a way Vulkan cannot reproduce.  If the shader cannot safely do
the write pattern, the Python Geo path cannot do it either.

Allowed alternatives must preserve source-owned execution.  Examples include:

\begin{itemize}[leftmargin=*]
    \item computing a source-owned response from source-owned contact data;
    \item using read-only frame-start snapshots of target state;
    \item accumulating source-owned diagnostic values;
    \item designing an explicit GLSL-compatible staged pass before adding the
    equivalent Python implementation.
\end{itemize}

Any proposal that needs one contact solve to write both particles must first
change the GLSL architecture, not just the Python model.

\subsection{Dynamics Versus Harness}

\texttt{GeoDynamics.py} should not:

\begin{itemize}
    \item read cfg files;
    \item create particle lists from cfg files;
    \item know about UI flags;
    \item know about reports or capture files;
    \item contain test harness iteration logic.
\end{itemize}

\texttt{GeoBase.py} may do those things because it is not intended to be a
one-to-one GLSL export unit.

\subsection{Current Standing Rules}

\begin{itemize}[leftmargin=*]
    \item Keep \texttt{GeoDynamics.py} one-to-one with GLSL export.
    \item Add contact handling only where it matches the shader model.
    \item Add reporting only through \texttt{Reporting.py}, not through
    dynamics.
    \item Record dynamics-rule changes in \texttt{ProblemEvolution.tex}.
\end{itemize}
